\begin{lemma}[The model density of \S4 inhabits (F1)--(F4) on $1<p<2$]
Let $n\ge1$ be an integer, let $1<p<2$, let $\varepsilon\ge0$ and let $a\ge0$.
Then the density
\[
f(x,z,h)=\bigl(\varepsilon+|h|^2\bigr)^{p/2}
\]
of \S4 --- which does not depend on $x$ or $z$ --- together with the modulus
$\sigma\equiv0$ satisfies
\[
\mathrm{VariationalStructureF123}\bigl(n,p,\varepsilon,\ \lambda=1,\
\Lambda=1,\ a,\ \nu_0=p(p-1),\ \nu_1=p,\ \Omega=\R^n,\ f,\ \sigma\bigr),
\]
and in addition (F4) holds for that $\sigma$ at every $c_\sigma>0$ and every
$\delta>0$:
\[
\sigma(t)\le c_\sigma t^\delta\qquad\text{for all }t\ge0 .
\]

\noindent\textbf{Why this node exists.}  The Note for the formalization requires
at l.~258 that a cited result ``must not be vacuous (each must be a true
statement for the operators to which it is applied)'', and the demand applies a
fortiori to the structure bundle on which every \S4 statement is conditioned: if
$\mathrm{VariationalStructureF123}$ were satisfied by no density, then
Theorem~Theorem~4 of the thesis, \noderef{MinimumGradientHolder} and everything proved from
them would be vacuously true.  This node discharges that demand for the density
\S4 is written about.  It is the (F1)--(F3) counterpart of the $\S\S1$--3
witness for (Q1) of the thesis--(Q2) of the thesis, and its standing use is the same: a
clause added to the bundle is re-checked against it.

\textbf{On the constants.}  They are admissible exactly where the consumers ask
for them.  $0<\lambda\le\Lambda$ holds as $1\le1$.  $0<\nu_0<\nu_1$ holds
because $0<p(p-1)<p$ precisely when $0<p-1<1$, i.e.\ on $1<p<2$; on $p\ge2$ the
same computation below gives the two constants in the other order, $\nu_0=p$ and
$\nu_1=p(p-1)$, and the statement is confined to the singular range for that
reason alone.  $a$ is left free and only required nonnegative, so the witness is
available at the consumers' $0<a$; the density in fact satisfies (F1) with
$a=0$, since the two bounds there hold with equality.  $\sigma\equiv0$ is
admissible because $f$ is independent of $x$ and $z$, so the left-hand side of
(F3) is identically $0$; a density with genuine $x$-dependence needs a genuine
$\sigma$, and nothing here claims otherwise.
\end{lemma}
\begin{proof}
Throughout write $w(h)=\varepsilon+|h|^2$ and $f(h)=w(h)^{p/2}$, suppressing the
constant dependence on $x$ and $z$.  Note $\varepsilon\ge0$ gives
$|h|^2\le w(h)$, and $1<p<2$ gives $p-2<0$ and $p-1\in(0,1)$.

\emph{Carath\'eodory, measurability.}  For fixed $z$ and $h$ the map
$x\mapsto f(h)$ is constant, hence Borel measurable.

\emph{Carath\'eodory, joint continuity.}  Fix $x$.  The map
$(z,h)\mapsto f(h)$ is the composition of the projection $(z,h)\mapsto h$ with
$f$, so it suffices that $f$ is continuous on $\R^n$.  If $\varepsilon>0$ then
$w\ge\varepsilon>0$ and $f$ is the composition of the smooth map $h\mapsto w(h)$
with $t\mapsto t^{p/2}$, which is smooth on $(0,\infty)$.  If $\varepsilon=0$
then $f(h)=|h|^{p}$, which is continuous on $\R^n\setminus\{0\}$ by the same
argument and continuous at the origin because $|h|^p\to0=f(0)$ as $h\to0$, using
$p>0$ and the real-power convention $0^{p}=0$ for $p\ne0$.

\emph{(F1).}  Both inequalities are immediate, with $\lambda=\Lambda=1$:
\[
-a+1\cdot w(h)^{p/2}\;\le\;f(h)\;=\;w(h)^{p/2}\;\le\;1\cdot w(h)^{p/2}+a ,
\]
since $a\ge0$.  This holds for every $x$, in particular for every $x\in\R^n$.

\emph{(F2), the $C^1$ clause.}  We claim $f\in C^1(\R^n)$.  For
$h\in D(\varepsilon)=\{w\ne0\}$ this is the chain rule as above, with
\begin{equation}
Df(h)\xi=p\,w(h)^{\frac p2-1}(h,\xi),
\tag{$\ast$}
\end{equation}
obtained from $D(|y|^2)\xi=2(h,\xi)$ at $y=h$ and $\frac p2\cdot2=p$.  If
$\varepsilon>0$ then $D(\varepsilon)=\R^n$ and we are done.  If
$\varepsilon=0$ the single point $h=0$ remains, where $f(h)=|h|^p$.  There
\[
\frac{|f(h)-f(0)-0|}{|h|}=\frac{|h|^p}{|h|}=|h|^{p-1}\xrightarrow[h\to0]{}0
\]
because $p-1>0$, so $f$ is differentiable at $0$ with $Df(0)=0$; and
$\|Df(h)\|=p|h|^{p-1}\to0=\|Df(0)\|$ by $(\ast)$ and Cauchy--Schwarz (the
supremum over $|\xi|\le1$ of $p|h|^{p-2}(h,\xi)$ is $p|h|^{p-1}$), so $Df$ is
continuous at $0$ as well.  Hence $f\in C^1(\R^n)$ in every case.

\emph{(F2), the $C^2$ clause.}  On the open set $D(\varepsilon)=\{w\ne0\}$ the
map $h\mapsto w(h)$ is a polynomial that does not vanish, so
$h\mapsto w(h)^{\frac p2-1}$ is $C^\infty$ there, and differentiating $(\ast)$
once more by the product and chain rules gives, for $h\in D(\varepsilon)$ and
$\zeta\in\R^n$,
\begin{equation}
D\bigl(h'\mapsto Df(h')\zeta\bigr)(h)\,\zeta
= p\,w(h)^{\frac p2-1}|\zeta|^2
  + p(p-2)\,w(h)^{\frac p2-2}(h,\zeta)^2 .
\tag{$\ast\ast$}
\end{equation}
In particular $f$ is $C^2$ on $D(\varepsilon)$, which is the clause as stated.

\emph{(F2), the two bounds.}  Fix $h\in D(\varepsilon)$ and $\zeta$, and write
$w=w(h)>0$.  By Cauchy--Schwarz and $|h|^2\le w$,
\[
0\;\le\;(h,\zeta)^2\;\le\;|h|^2|\zeta|^2\;\le\;w\,|\zeta|^2 .
\]
Since $p(p-2)<0$, the second term of $(\ast\ast)$ is nonpositive, which gives
the upper bound at once:
\[
D\bigl(h'\mapsto Df(h')\zeta\bigr)(h)\,\zeta\;\le\;p\,w^{\frac p2-1}|\zeta|^2
=\nu_1\,w^{\frac p2-1}|\zeta|^2 .
\]
For the lower bound, replacing $(h,\zeta)^2$ by the larger $w|\zeta|^2$ in a
term with the negative factor $p(p-2)$ decreases the sum, so using
$w^{\frac p2-2}\cdot w=w^{\frac p2-1}$ (valid since $w>0$),
\[
D\bigl(h'\mapsto Df(h')\zeta\bigr)(h)\,\zeta
\;\ge\;p\,w^{\frac p2-1}|\zeta|^2+p(p-2)\,w^{\frac p2-1}|\zeta|^2
=p(p-1)\,w^{\frac p2-1}|\zeta|^2=\nu_0\,w^{\frac p2-1}|\zeta|^2 .
\]
Both are the required bounds, and both hold for every $x$.

\emph{The modulus clauses.}  $\sigma\equiv0$ is bounded above by $C=0$, is
continuous, is monotone, and satisfies $\sigma(0)=0$.

\emph{(F3).}  For all $x,x'$, all $z,z'$ and all $h$ the left-hand side is
$|f(h)-f(h)|=0$, because $f$ does not depend on $x$ or $z$, and the right-hand
side is $(1+|h|^2)^{p/2}\cdot\sigma(\cdot)=0$ as well.  So (F3) holds with
equality.

\emph{(F4).}  For $t\ge0$ and $c_\sigma>0$, $\delta>0$ we have
$\sigma(t)=0\le c_\sigma t^\delta$, since $t^\delta\ge0$ for $t\ge0$ under the
real-power convention.  This is the second conjunct.
\end{proof}
